\section{Study Design}
\label{sec:design}

\subsection{Research questions}
\label{sec:rqs}
The study has two parts.
The first measures what supplying designated evidence does to GPT predictions, both in aggregate and claim by claim, which exposes the regressions that aggregate accuracy hides.
The second describes those regressions through blinded judge panels whose readings of the same evidence can be compared across model families.
This study seeks to answer four questions, two for each part:
% COAUTHOR REVIEW: RQ3 and RQ4 are reworded from "diagnostic failure modes" to descriptive adjudication outcomes.
\begin{description}
\item[RQ1] How does providing FEVER-designated gold evidence affect aggregate LLM fact-verification accuracy?
\item[RQ2] When evidence changes a previously correct prediction into an incorrect one, how frequent and consistent are those regressions across models?
\item[RQ3] Which observable adjudication outcomes---decisiveness, agreement with FEVER, and the disclosure stage at which they change---characterize evidence-induced regressions under progressive disclosure?
\item[RQ4] Which of those outcomes depend on the adjudicating model family?
\end{description}

\subsection{FEVER dataset and sampling}
We sample from the FEVER development set \citep{thorne-etal-2018-fever}.
The manuscript experiment is a balanced 10{,}000-claim extract: 5{,}000 Supported and 5{,}000 Refuted, seed 42, \texttt{NOT ENOUGH INFO} excluded.
Gold evidence is resolved from official FEVER Wikipedia shards into the text shown to the models.
Sampling and validation artifacts are committed with the frozen repository; third-party FEVER and Wikipedia text are not relicensed by the project MIT license.

\subsection{Models and paired evidence conditions}
Each claim is evaluated four times: GPT-5.4 and GPT-5.4-mini, each under \texttt{claim\_only} and \texttt{claim\_plus\_evidence}.
That yields 40{,}000 predictions.
The evidence condition places the designated FEVER evidence (gold evidence) directly in the prompt; no retrieval step is involved.
Because the two conditions differ only in whether that evidence is present, a change in a claim's prediction between them reflects the supplied evidence rather than retrieval.
Original GPT runs stored predicted labels only, so later analysis cannot recover confidence, rationales, or token-level attributions.

\subsection{Transition definitions}
For each model, a claim has one of four paired transitions:
\emph{robust} (correct without and with evidence),
\emph{rescue} (wrong without, correct with),
\emph{resistant} (wrong in both),
or \emph{regression} (correct without, wrong with).
A unique-claim regression is a claim that regresses for at least one of the two GPT models.
Unique-claim counts are never treated as model-observations: 115 GPT-5.4 regressions and 151 mini regressions are 226 unique claims, not 266.

\subsection{Analysis v2}
Analysis v2 treats the 40{,}000-row experiment file as immutable source data.
It rebuilds the paired dataset, runs McNemar tests and bootstrap confidence intervals, extracts linguistic and evidence-structure features, embeds claim--evidence pairs, and applies two local NLI models (nli-deberta-v3-small and distilbert-base-uncased-mnli) as diagnostics, not as ground truth.
Leakage audits and an independent 69-check headline verifier are part of the frozen record.
Confirmatory tests were pre-specified; NLI and feature analyses are exploratory.

\subsection{Diagnostic cohort}
The adjudication cohort contains 1{,}061 unique claims: all 226 unique regressions, 335 resistant, 250 rescue, and 250 robust.
One item (\texttt{SA-000352}) is provider-filtered, so judged $N=1{,}060$.
The Grok panel judges the 1{,}060-claim cohort; the Claude panel judges only the 226 regressions.
This diagnostic sample is not representative of all 10{,}000 claims, and cross-family comparisons always pair the same 226 claims.

\subsection{Progressive-disclosure adjudication}
\label{sec:protocol}
Each judge is a single isolated run of a locked model that applies a fixed rubric to each claim and its evidence and returns a verdict (Supported, Refuted or Ambiguous) with a rating of evidence sufficiency; code combines the five judges' outputs into a consensus that is Supported, Refuted, Ambiguous or Unresolved.
The design discloses the evidence in cumulative stages so that a panel's verdicts show whether a claim can be decided from the sentences GPT saw, is decided only after more context, or remains undecided.
Five isolated judges per stage and batch see:
Stage~A, the original selected gold-evidence sentence text (identical to the GPT evidence prompt);
Stage~B, Stage~A plus the page titles;
Stage~C, structured FEVER evidence---ordered sentence pointers, annotation boundaries, and the text of available alternative evidence sets.
Stage~C can therefore add information as well as change representation: on the repaired inputs it adds sentence text beyond the selected sentences for 149 cohort claims, including 41 regressions.
Judges do not browse, retrieve new pages, or see FEVER labels, GPT predictions, cohort tags, NLI scores, or other judges.
Each stage is frozen before the next stage's packets are built, and packets carry only current-stage evidence.
Consensus is computed by code from the five raw judgments; resolver outputs never replace raw consensus.
The Grok panel is locked to \texttt{cursor-grok-4.6-high-fast} and the Claude panel to \texttt{claude-opus-5-thinking-high}.

\subsection{Descriptive taxonomy}
\label{sec:taxonomy}
For each claim the taxonomy records decisiveness (Supported or Refuted versus Ambiguous or Unresolved) and FEVER agreement at every stage, the first gold-agreeing stage, title and Stage~C resolutions, reversals, and losses of decisiveness.
Each claim then receives exactly one category, in priority order:
\emph{final nondecisive} (Stage~C consensus Ambiguous or Unresolved);
\emph{decisive judge--FEVER disagreement} (decisive Stage~C verdict opposite to FEVER);
\emph{nonmonotonic judge path} (final gold agreement after a reversal or loss of decisiveness);
\emph{sentence-only gold-agreeing} (gold-agreeing from Stage~A onward);
\emph{title-disclosure gold-agreeing} (first gold agreement at Stage~B);
and, for first gold agreement at Stage~C, \emph{additional-evidence gold-agreeing} when Stage~C adds sentence text and \emph{structured-disclosure gold-agreeing} otherwise.
These categories describe where a blinded panel's verdict stands; they are not causal mechanisms inside GPT.
A decisive disagreement with FEVER is not evidence that the FEVER label is wrong, and sentence-only gold agreement is compatible with, but does not establish, a failure by GPT to use the supplied evidence.
The taxonomy was specified post hoc, after inspection of the historical results, and awaits coauthor review.

\subsection{Evidence audit and Stage~C correction}
\label{sec:correction}
The historical Stage~C packets took each evidence line's last tab-separated field, often a hyperlink target, instead of the sentence text.
Chosen-set text differed from the canonical sentences for 936 of 1{,}061 cohort claims (190 of the 226 regressions); ten further claims lacked structured text because decomposed-Unicode titles were missing from a cache.
Exact selected evidence sets were recovered by recomputing the original set identifiers, and all selected and alternative sentences were rebuilt from local archival FEVER data.
Stage~A inputs equal the GPT evidence prompts and Stage~B inputs equal the historical records exactly; their frozen judgments are retained under the documented fresh-context-per-stage protocol, and all Stage~C judgments are new.
Consensus paths therefore combine historical Stage~A/B judgments with new Stage~C judgments.

A post hoc correction protocol, frozen on 2026-09-27 before any new judgment, fixed the repaired packets, the taxonomy, the statistical analyses, and the execution rules, including acceptance of the first schema-valid response for each job.
Seventy required jobs (55 Grok, 15 Claude; 6{,}430 new Stage~C judgments) all delivered schema-valid outputs.

\paragraph{Post-judgment amendment.}
Execution failures meant that some jobs needed more attempts than the protocol allowed and that its execution-record rules could not be met as written.
Retries never depended on model outputs.
After all outputs existed and had been inspected, we adopted a dated amendment (2026-09-28) that relaxes those execution rules and specifies that each job's analysed output is the output its execution procedure accepted.
% COAUTHOR REVIEW: accepted outputs are the main selection by the authors' post-judgment decision; first-valid is sensitivity.
For three jobs (two Claude, one Grok) an earlier schema-valid output exists, so we repeat every analysis with those first valid outputs substituted (Appendix~\ref{app:amendment}).
The amendment is not preregistered; under the original rules the protocol's primary analysis remains incomplete.
Packets, models, taxonomy, statistics and analysis code are unchanged, and the protocol's follow-up panels (a newly selected residual Claude panel and a Stage~C resolver) were not run.
Execution records and provenance evidence are documented in the released artifact.

\subsection{Supporting cross-family data}
\label{sec:cross-family-design}
GLM-5.3 contributes a completed Stage~A panel compared to frozen Cursor Stage~A ($n=1{,}060$); Stage~A inputs are unaffected by the audit.
A historical Claude Stage~C test on 231 Grok leftovers used the defective Stage~C packets and a residual set selected with defective Grok Stage~C judgments; it is reported only as a historical observation (Appendix~\ref{app:historical}).

\subsection{Statistical analysis}
Paired accuracy uses McNemar tests; regression and rescue rates use bootstrap confidence intervals.
Descriptive adjudication rates use Wilson 95\% intervals.
Cross-family comparisons on the paired 226 claims use exact McNemar tests with Benjamini--Hochberg correction across 14 tests; shared-versus-model-specific nondecisiveness uses odds ratios in a separate two-test family.
Cross-family taxonomy agreement is reported as exact match and Cohen's $\kappa$.
All adjudication analyses are post hoc.
